\documentclass{article}
\usepackage{amsmath,amssymb}
\usepackage{xurl}
\usepackage[hidelinks]{hyperref}
% Shared commands for notes in docs/paper-gaps/.

\DeclareUnicodeCharacter{2080}{\ifmmode{}_{0}\else\textsubscript{0}\fi}
\DeclareUnicodeCharacter{2081}{\ifmmode{}_{1}\else\textsubscript{1}\fi}
\DeclareUnicodeCharacter{2082}{\ifmmode{}_{2}\else\textsubscript{2}\fi}
\DeclareUnicodeCharacter{2099}{\ifmmode{}_{n}\else\textsubscript{n}\fi}

\newcommand{\Lean}{\textsc{Lean}}
\ExplSyntaxOn
\NewDocumentCommand{\leanid}{m}{
  \ifmmode
    \textsf{\detokenize{#1}}
  \else
    \group_begin:
    \urlstyle{sf}
    \tl_set:Nn \l_tmpa_tl {#1}
    \tl_replace_all:Nnn \l_tmpa_tl {\_} {_}
    \exp_args:NV \path \l_tmpa_tl
    \group_end:
  \fi
}
\ExplSyntaxOff
\providecommand{\tr}{\operatorname{tr}}
\newcommand{\tnleanIssueBase}{https://github.com/LionSR/TNLean/issues/}
\newcommand{\tnleanPrBase}{https://github.com/LionSR/TNLean/pull/}
\newcommand{\tnleanDocBase}{https://sirui-lu.com/TNLean/docs/}
\newcommand{\ghissue}[1]{\href{\tnleanIssueBase#1}{GitHub issue \##1}}
\newcommand{\ghpr}[1]{\href{\tnleanPrBase#1}{PR \##1}}
\newcommand{\doclink}[1]{\href{\tnleanDocBase#1.html}{\path{#1}}}

% Dirac notation and the identity operator, as in blueprint/src/macros/common.tex.
% \providecommand keeps note-local definitions authoritative.
\usepackage{dsfont}
\providecommand{\ket}[1]{|#1\rangle}
\providecommand{\bra}[1]{\langle #1|}
\providecommand{\braket}[2]{\langle #1 | #2 \rangle}
\providecommand{\ketbra}[2]{|#1\rangle\!\langle #2|}
\providecommand{\Id}{\mathds{1}}
\providecommand{\one}{\mathds{1}}
\providecommand{\C}{\mathbb{C}}
\providecommand{\MN}[1]{M_{#1}(\C)}
\providecommand{\MD}{M_D(\C)}

% External referencing for paper sources.  The build helper writes sanitized
% auxiliary files under build/paper-aux/ and excludes duplicated source labels.
\usepackage{xr}

% Import a generated paper auxiliary file with a caller-supplied label prefix.
% The candidate paths support builds from docs/paper-gaps/, the repository root,
% and build/paper-aux/ itself.
\newcommand{\importpaperaux}[2]{%
  \IfFileExists{../../build/paper-aux/#2.aux}{%
    \externaldocument[#1]{../../build/paper-aux/#2}%
  }{%
    \IfFileExists{build/paper-aux/#2.aux}{%
      \externaldocument[#1]{build/paper-aux/#2}%
    }{%
      \IfFileExists{#2.aux}{%
        \externaldocument[#1]{#2}%
      }{}%
    }%
  }%
}

% General external reference macros
\newcommand{\paperref}[2]{\ref{#1-#2}}
\newcommand{\papereqref}[2]{(\ref{#1-#2})}

% CPSV16 source (1606.00608 / MPDO-22-12-17-2.tex) external document and shortcuts
\importpaperaux{cpsv16-}{1606.00608/MPDO-22-12-17-2-tnlean}
\newcommand{\cpsvref}[1]{\paperref{cpsv16}{#1}}
\newcommand{\cpsveqref}[1]{\papereqref{cpsv16}{#1}}

% Verdict marker: \gapnote{<kind>}{<status>}. Kind is the policy
% classification (clarification, local-correction, scope-restriction,
% unfaithful, false-source, open-gap); status is open, wip (elimination
% underway), resolved, or historical. Machine-read by the site index;
% prints a verdict line.
\newcommand{\gapnote}[2]{\par\noindent\textbf{Verdict:} #1 (#2).\par}

\title{The Supplied Fixed Pair in the Reduced-Representative Step of Proposition IV.5}
\author{}
\date{2026-08-16}
\begin{document}
\maketitle

\section{Source statement and derivation of the fixed pair}
Proposition IV.5 of Cirac--P\'erez-Garc\'ia--Schuch--Verstraete states that
for a continuous map of tensors $[0,1]\ni x\mapsto\mathcal W(x)$, each
generating MPUs but not necessarily in canonical form, the ranks $r(x)$ and
$\ell(x)$ of the associated simple tensors are constant
\cite[Proposition~IV.5]{Cirac2017MPU}. The proof begins with a single tensor
$\mathcal W$ generating an MPU and derives, from the MPU property alone and
after blocking $D^4$ sites, that the blocked canonical form is simple and
that its transfer operator is rank one: as a completely positive map it is
\begin{align}
  X&\longmapsto\tr(LX)\,R,
\end{align}
with positive semidefinite left and right fixed points $L\succeq0$ and
$R\succeq0$ normalized by $\tr(LR)=1$. The proof then reads off, from the
displayed factorization through an isometry $V$, that the tensor absorbs the
support projections on both sides,
\begin{align}
  \mathcal WP&=\mathcal W=Q\mathcal W,
\end{align}
where $P$ and $Q$ are the orthogonal projections onto the supports of $L$ and
$R$, respectively. The source treats all four ingredients---the rank-one
transfer formula, the positive semidefiniteness, the normalization, and the
two-sided support absorption---as derived consequences of blocking an
arbitrary MPU tensor.
% Source lines: paper_v2.tex 747--756. The isometry $V$ is the one displayed in
% the figure IV_index1.png referenced at paper line 754.

\section{The formal scope issue}
The formal reduced-representative calculation reverses the direction of this
part of the argument: it takes the fixed pair and its consequences as explicit
hypotheses rather than deriving them. The letterwise absorptions
$W^{ij}P=W^{ij}$ and $QW^{ij}=W^{ij}$, the rank-one transfer formula
$E_W(X)=\tr(LX)R$, the positivity $L,R\succeq0$, and the normalization
$\tr(LR)=1$ are all assumptions of the formal statements. No hypothesis is
inferred from the bare MPU predicate; the module docstring records this
candor in the sentence that all fixed-point and support hypotheses are
explicit.

What the formal results prove is the remainder of the source argument,
starting from the supplied pair. With $T=\widetilde P$ the reduced projection
of $P$ and $Q$, the reduced tensor $\widetilde W^{ij}=TW^{ij}T$ generates the
same periodic MPO at every system size, its transfer operator has the
compressed rank-one formula with left and right fixed points $TLT$ and $TRT$,
the compression preserves the trace normalization
$\tr((TLT)(TRT))=1$, and in isometric coordinates $V^\dagger V=I$
with $VV^\dagger=T$, both compressed fixed matrices $V^\dagger LV$ and
$V^\dagger RV$ are positive definite.

The restriction is not vacuous. The source derivation of the fixed pair uses
the blocking to $D^4$ sites and the simple canonical-form structure of the
blocked tensor; the unblocked tensor, or an arbitrary representative of it,
carries no such pair without that construction. Until the supplier is
formalized, the formal statements apply only to tensors already equipped with
the fixed pair, whereas the source claims the conclusions for every MPU
tensor after blocking.

\section{Affected declarations}
The declaration
\leanid{MPOTensor.mpo\_virtualSandwich\_reducedProjection\_eq} assumes the
letterwise absorptions $W^{ij}P=W^{ij}$ and $QW^{ij}=W^{ij}$.
\leanid{MPOTensor.transferMap\_virtualSandwich\_reducedProjection} assumes the
rank-one transfer formula. The declarations
\leanid{MPOTensor.trace\_compressed\_fixed\_pair\_reducedProjection} and
\leanid{MPOTensor.compressed\_fixed\_pair\_posDef\_reducedProjection} assume,
respectively, the two-sided support absorptions of $L$ and $R$ together with
$\tr(LR)=1$, and the positive pair $L,R\succeq0$ with its support projections;
the latter also assumes an isometry $V$ onto the range of $T$.

The comparison declaration
\leanid{MPOTensor.virtualSandwich\_hat\_eq\_reducedProjection} assumes the
factor equations $XL=P$, $RZ=Q$, and $PQY=T$ together with the letterwise
absorptions. The rank equalities
\leanid{MPOTensor.rightRank\_hat\_eq\_reducedProjection} and
\leanid{MPOTensor.leftRank\_hat\_eq\_reducedProjection} use the same supplied
data and additionally assume that $X$, $Z$, and $Y$ are invertible. Finally,
\leanid{MPOTensor.exists\_units\_for\_hat\_reducedProjection} assumes positive
semidefinite fixed matrices $L$ and $R$ and supplies the invertible ambient
support-inverse extensions $X$ and $Z$, together with an invertible
right factor $Y$ satisfying $PQY=T$.

\section{Elimination plan}
The restriction is removed by formalizing the missing supplier: the blocking
to $D^4$ sites of an arbitrary MPU tensor, the two-projection Jordan
decomposition of the resulting support projections, and the isometric
factorization that yields the letterwise absorptions and the rank-one
transfer formula with positive semidefinite normalized fixed points. This is
tracked by issues~\#6137 and~\#6022: the former constructs and transports the
reduced representative and supplies the positive fixed pair with the
letterwise absorptions for the common blocked tensor, and the latter carries
the full continuity-index route of which Proposition~IV.5 is a part. Until
these are closed, no formal statement may infer the fixed pair or the
absorptions from the bare MPU predicate.

\section{Verdict}
\textbf{Verdict: scope restriction.} The paper derives the fixed pair
$L,R\succeq0$ with $\tr(LR)=1$, the rank-one transfer formula, and the
two-sided support absorption $\mathcal WP=\mathcal W=Q\mathcal W$ from the
MPU property after blocking $D^4$ sites. The present formalization states the
reduced-representative conclusions under these ingredients as explicit
hypotheses and does not construct the supplier. The results are correct as
stated, but narrower than the source argument; the hypotheses become
theorems only after issues~\#6137 and~\#6022 are resolved.

\bibliographystyle{alpha}
\bibliography{references}
\end{document}